\documentclass{article}
\usepackage[colour]{imprimatur}
\begin{document}

\section{Rectangles}

The area of a rectangle is the product of its sides; swapping the sides
keeps it \leancheck[stmt=ef984ed62ab8, sig=f639e6d6211a]{Shapes.Rect.area_swap, Shapes.Rect.area}.

\begin{leanblock}[sum=v80Pqxg1yV]{Shapes.Rect.area_swap}{Swap}
\LeanCode{%
\LeanKeyword{theorem} \LeanConst{Rect.area\_swap} (\LeanBound{r} : \LeanConst{Rect}) : \{ w := \LeanBound{r}.h, h := \LeanBound{r}.w \}.area = \LeanBound{r}.area}
\end{leanblock}

As mathematics:
\begin{leanblock}[view=math, sum=rtzvrpiUKQ]{Shapes.Rect.area_swap}{Swap}
\LeanMath{|r_{\mathrm{h}} \times r_{\mathrm{w}}| = |r|}
\end{leanblock}

The host's own executable, a module named \texttt{Main} like the
extractor's, is the one the paper sees:
\begin{leanblock}[view=lean, sum=7Df85bEOvD]{main}{Entry point}
\LeanCode{%
\LeanKeyword{def} \LeanConst{main} : \LeanConst{IO} \LeanConst{Unit}}
\end{leanblock}

\end{document}
